\chapter[Single-sorted presentation]{Single-sorted presentation of higher-order categories}
\label{ch:single_sorted}

In the single-sorted presentation of higher-order categories, objects, morphisms, 2-morphisms, and
all higher morphisms live in a single set\footnote{%
  Throughout this text, following the convention of~\cite{vidal2024higher}, we use the word
  \emph{set} in the classical mathematical sense. However, the reader should be aware that, in
  general, the accompanying Lean~4 formalization works with types rather than the Lean concept of
  sets.%
}, distinguished by source and target operations at each dimension. Also, composition at each
dimension is defined as a partial operation that is only defined when appropriate source-target
matching conditions are satisfied. In this chapter, we present the very basic definitions and axioms
of single-sorted categories and single-sorted functors between them.

\begin{notation}
  Throughout the text, if not otherwise specified, we will use $\Index$ to denote a set of indices
  equipped with a preorder relation $<$. This general setting is used to define the core structures.
  In practice, the main definitions specialize $\Index$ to either a finite set $n = \set{0, 1,
  \ldots, n - 1}$ or the set $\omega$ of natural numbers.
\end{notation}

\section{Single-sorted categories}

We start by defining the necessary data for defining a single-sorted category.

\begin{definition}[Single-sorted category structure]
  \label{def:SingleSortedCategoryStruct}
  \lean{HigherCategoryTheory.SingleSorted.CategoryStruct}
  \leanok
  Let $C$ be a set whose elements will be called \emph{morphisms}. The required data for defining a
  single-sorted $\Index$-category structure on $C$ consists of:

  \begin{enumerate}
    \item \label{ax:ss_sc} A unary operation $\Sc^k \colon C \to C$ for each index $k \in \Index$,
      called the \emph{source} at dimension $k$.
    \item \label{ax:ss_tg} A unary operation $\Tg^k \colon C \to C$ for each index $k \in \Index$,
      called the \emph{target} at dimension $k$.
    \item \label{ax:ss_comp} A partial binary operation $\Pcomp^k \colon C \times C \rightharpoonup
      C$ for each index $k \in \Index$, called the \emph{composition} at dimension $k$.
  \end{enumerate}

  The composition $g \Pcomp^k f$ is defined if and only if $\Sc^k (g) = \Tg^k (f)$, in which case we
  say that $g$ and $f$ are \emph{composable} at dimension $k$.

  We will represent the above data as a tuple $\cat{C} = (C, (\Sc^k, \Tg^k, \Pcomp^k)_{k \in
  \Index})$.
\end{definition}

We split the axioms of single-sorted categories into two parts: first, those concerning only one
dimension at a time, and then those concerning the interaction between different dimensions. This
division is made by first defining the notion of pre-single-sorted category:

\begin{definition}[Pre-single-sorted category]
  \label{def:PreSingleSortedCategory}
  \uses{def:SingleSortedCategoryStruct}
  \lean{HigherCategoryTheory.SingleSorted.PreCategory}
  \leanok
  A \emph{pre-single-sorted $\Index$-category} is a tuple
  \[
    \cat{C} = (C, (\Sc^k, \Tg^k, \Pcomp^k)_{k \in \Index}),
  \]
  satisfying, for each index $k \in \Index$, the following axioms:

  \begin{enumerate}
    \item \label{ax:ss_glob} For all $f \in C$,
      \begin{align*}
        \Sc^k (\Sc^k (f)) = \Sc^k (f), \qquad & \Tg^k (\Sc^k (f)) = \Sc^k (f), \\
        \Sc^k (\Tg^k (f)) = \Tg^k (f), \qquad & \Tg^k (\Tg^k (f)) = \Tg^k (f).
      \end{align*}
    \item \label{ax:ss_sc_tg_comp} For all $f, g \in C$, if $g$ and $f$ are composable, then
      \begin{align*}
        \Sc^k (g \Pcomp^k f) = \Sc^k (f), \qquad & \Tg^k (g \Pcomp^k f) = \Tg^k (g).
      \end{align*}
    \item \label{ax:ss_right_id} For all $f \in C$, the morphisms $f$ and $\Sc^k (f)$ are composable
      and
      \[
        f \Pcomp^k \Sc^k (f) = f.
      \]
    \item \label{ax:ss_left_id} For all $f \in C$, the morphisms $\Tg^k (f)$ and $f$ are composable
      and
      \[
        \Tg^k (f) \Pcomp^k f = f.
      \]
    \item \label{ax:ss_assoc} For all $f, g, h \in C$, if $h$ and $g$ are composable, and $g$ and
      $f$ are composable, then $h \Pcomp^k g$ and $f$ and $h$ and $g \Pcomp^k f$ are composable, and
      \[
        (h \Pcomp^k g) \Pcomp^k f = h \Pcomp^k (g \Pcomp^k f).
      \]
  \end{enumerate}
\end{definition}

We finally define a single-sorted category structure by extending the above definition with the
axioms concerning the interaction between different dimensions:

\begin{definition}[Single-sorted category]
  \label{def:SingleSortedCategory}
  \uses{def:PreSingleSortedCategory}
  \lean{HigherCategoryTheory.SingleSorted.Category}
  \leanok
  A \emph{single-sorted $\Index$-category} is a pre-single-sorted $\Index$-category $\cat{C} = (C,
  (\Sc^k, \Tg^k, \Pcomp^k)_{k \in \Index})$ satisfying, in addition to the axioms of
  pre-single-sorted categories, the following cross-dimensional axioms, for all indices $k, j \in
  \Index$ with $j < k$:

  \begin{enumerate}
    \item \label{ax:ss_cross_glob} For all $f \in C$,
      \begin{align*}
        \Sc^k (\Sc^j (f)) = \Sc^j (f), \qquad & \Sc^j (\Sc^k (f)) = \Sc^j (f), \\
        \Sc^j (\Tg^k (f)) = \Sc^j (f), \qquad & \Tg^k (\Tg^j (f)) = \Tg^j (f), \\
        \Tg^j (\Tg^k (f)) = \Tg^j (f), \qquad & \Tg^j (\Sc^k (f)) = \Tg^j (f).
      \end{align*}
    \item \label{ax:ss_cross_sc_tg_comp} For all $f, g \in C$, if $g$ and $f$ are composable at
      dimension $j$, then $\Sc^k (g)$ and $\Sc^k (f)$ and $\Tg^k (g)$ and $\Tg^k (f)$ are also
      composable at dimension $j$, and
      \begin{align*}
        \Sc^k (g \Pcomp^j f) = \Sc^k (g) \Pcomp^j \Sc^k (f), \\
        \Tg^k (g \Pcomp^j f) = \Tg^k (g) \Pcomp^j \Tg^k (f).
      \end{align*}
    \item \label{ax:ss_interchange} For all $f_1, f_2, g_1, g_2 \in C$, suppose that:
      \begin{itemize}
        \item $g_2$ and $f_2$ are composable at dimension $j$,
        \item $g_1$ and $f_1$ are composable at dimension $j$,
        \item $g_2$ and $g_1$ are composable at dimension $k$, and
        \item $f_2$ and $f_1$ are composable at dimension $k$.
      \end{itemize}
      Then, the morphisms $(g_2 \Pcomp^j f_2)$ and $(g_1 \Pcomp^j f_1)$ are composable at dimension
      $k$, and the morphisms $(g_2 \Pcomp^k g_1)$ and $(f_2 \Pcomp^k f_1)$ are composable at
      dimension $j$, and
      \[
        (g_2 \Pcomp^j f_2) \Pcomp^k (g_1 \Pcomp^j f_1) = (g_2 \Pcomp^k g_1) \Pcomp^j (f_2
        \Pcomp^k f_1).
      \]
      That is, composing first at dimension $j$ and then at dimension $k$ yields the same result as
      composing first at dimension $k$ and then at dimension $j$.
  \end{enumerate}
\end{definition}

\begin{definition}[Single-sorted finite category]
  \label{def:SingleSortedNCategory}
  \uses{def:SingleSortedCategory}
  \lean{HigherCategoryTheory.SingleSorted.NCategory}
  \leanok
  Given a natural number $n \in \NN$, a \emph{single-sorted $n$-category} is a single-sorted
  category indexed by the finite set $n = \set{0, 1, \ldots, n-1}$.
\end{definition}

\begin{remark}
  \label{thm:PreCategoryLift}
  \uses{def:SingleSortedNCategory}
  \lean{HigherCategoryTheory.SingleSorted.PreCategory.lift}
  \leanok
  Any pre-single-sorted $1$-category is, in fact, a single-sorted $1$-category. Since the index
  set $1 = \set{0}$ has exactly one element, there are no pairs of distinct indices $j < k$, so all
  cross-dimensional axioms are vacuously satisfied.
\end{remark}

We now define the notion of cell or identity morphism at a given dimension in a single-sorted
category.

\begin{definition}[Cell]
  \label{def:cell}
  \uses{def:SingleSortedCategory}
  \lean{
    HigherCategoryTheory.SingleSorted.cell,
    HigherCategoryTheory.SingleSorted.cell_sc,
    HigherCategoryTheory.SingleSorted.cells,
    HigherCategoryTheory.SingleSorted.cells_sc
  }
  \leanok
  Let $\cat{C}$ be a single-sorted $\Index$-category and $k \in \Index$. We say that a morphism $f
  \in C$ is a \emph{$k$-cell} or \emph{identity morphism at dimension $k$} if and only if
  \[
    \Sc^k (f) = f.
  \]

  We will denote the set of $k$-cells of $\cat{C}$ by $C_k$.
\end{definition}

The above definition defines $k$-cells using the source operation, however, it is also possible to
define them using the targer operation and, using the axioms of single-sorted categories, one can
prove that both definitions are equivalent.

\begin{definition}[Target-based cell]
  \label{def:cell_tg}
  \uses{def:SingleSortedCategory}
  \lean{HigherCategoryTheory.SingleSorted.cell_tg, HigherCategoryTheory.SingleSorted.cells_tg}
  \leanok
  Let $\cat{C}$ be a single-sorted $\Index$-category and $k \in \Index$. A morphism $f \in C$ is a
  \emph{target-based $k$-cell} if and only if
  \[
    \Tg^k (f) = f.
  \]
\end{definition}

\begin{theorem}[Equivalence of cell definitions]
  \label{thm:cell_sc_iff_cell_tg}
  \uses{def:cell, def:cell_tg}
  \lean{
    HigherCategoryTheory.SingleSorted.cell_sc_iff_cell_tg,
    HigherCategoryTheory.SingleSorted.cells_sc_eq_cells_tg
  }
  \leanok
  Let $\cat{C}$ be a single-sorted $\Index$-category and $k \in \Index$. A morphism $f \in C$ is a
  $k$-cell if and only if it is a target-based $k$-cell.

  It follows that the set of $k$-cells $C_k$ is equal to the set of target-based $k$-cells.
\end{theorem}

\begin{proof}
  \leanok
  Let $f \in C$. Suppose first that $f$ is a $k$-cell, i.e., $\Sc^k (f) = f$. By
  axiom~\refaxiom{def:PreSingleSortedCategory}{ax:ss_glob}, we have that
  \[
    \Tg^k (f) = \Tg^k (\Sc^k (f)) = \Sc^k (f) = f,
  \]
  so $f$ is a target-based $k$-cell.

  Conversely, suppose that $f$ is a target-based $k$-cell, i.e., $\Tg^k (f) = f$. By
  axiom~\refaxiom{def:PreSingleSortedCategory}{ax:ss_glob}, we have that
  \[
    \Sc^k (f) = \Sc^k (\Tg^k (f)) = \Tg^k (f) = f,
  \]
  so $f$ is a $k$-cell.
\end{proof}

\begin{definition}[Single-sorted omega category]
  \label{def:SingleSortedOmegaCategory}
  \uses{def:cell}
  \lean{HigherCategoryTheory.SingleSorted.OmegaCategory}
  \leanok
  A \emph{single-sorted $\omega$-category} is a single-sorted category indexed by the set of natural
  numbers $\NN$, satisfying and additional axiom: for all $f \in C$, there exists a dimension $k \in
  \NN$ such that $f$ is a $k$-cell.
\end{definition}

\section{Single-sorted functors}

Unlike traditional category theory where functors act separately on objects and morphisms, here a
single-sorted functor is simply a function on the set of morphisms that preserves sources, targets,
and composition at each dimension.

\begin{definition}[Single-sorted functor]
  \label{def:SingleSortedFunctor}
  \uses{def:SingleSortedCategory}
  \lean{HigherCategoryTheory.SingleSorted.Functor}
  \leanok
  Let $\cat{C} = (C, (\Sc^k, \Tg^k, \Pcomp^k)_{k \in \Index})$ and $\cat{D} = (D, (\Sc^k, \Tg^k,
  \Pcomp^k)_{k \in \Index})$ be single-sorted $\Index$-categories. A \emph{single-sorted
  $\Index$-functor} from $\cat{C}$ to $\cat{D}$ is a map $F \colon C \to D$ satisfying, for each
  index $k \in \Index$, the following preservation properties:

  \begin{enumerate}
    \item \label{ax:ss_fun_sc_tg} For all $f \in C$,
      \[
        F (\Sc^k (f)) = \Sc^k (F (f)), \qquad F (\Tg^k (f)) = \Tg^k (F (f)).
      \]
    \item \label{ax:ss_fun_comp} For all $f, g \in C$, if $g$ and $f$ are composable at dimension
      $k$, then $F (g)$ and $F (f)$ are also composable at dimension $k$, and
      \[
        F (g \Pcomp^k f) = F (g) \Pcomp^k F (f).
      \]
  \end{enumerate}
\end{definition}

We can define composition of single-sorted functors as the usual composition of functions, and the
identity single-sorted functor as the identity function on the set of morphisms. The following
results assert that these constructions yield valid single-sorted functors.

\begin{theorem}[Composition of single-sorted functors]
  \label{thm:SingleSortedFunctor_comp}
  \uses{def:SingleSortedFunctor}
  \lean{HigherCategoryTheory.SingleSorted.Functor.comp}
  \leanok
  Let $\cat{C}, \cat{D}, \cat{E}$ be single-sorted $\Index$-categories, and let $F \colon \cat{C}
  \to \cat{D}$ and $G \colon \cat{D} \to \cat{E}$ be single-sorted $\Index$-functors. Then, the
  composition $G \circ F \colon \cat{C} \to \cat{E}$ is also a single-sorted $\Index$-functor.
\end{theorem}

\begin{proof}
  \leanok
  Let $k \in \Index$. For all $f \in C$, since $F$ and $G$ are single-sorted functors, we have that
  \[
    (G \circ F) (\Sc^k (f)) = G (F (\Sc^k (f))) = G (\Sc^k (F (f)))
    = \Sc^k (G (F (f))) = \Sc^k ((G \circ F) (f)),
  \]
  and similarly,
  \[
    (G \circ F) (\Tg^k (f)) = G (F (\Tg^k (f))) = G (\Tg^k (F (f)))
    = \Tg^k (G (F (f))) = \Tg^k ((G \circ F) (f)),
  \]
  that is, $G \circ F$ preserves sources and targets at dimension $k$.

  Now, let $f, g \in C$ be composable at dimension $k$. We have that
  \[
    \begin{split}
      (G \circ F) (g \Pcomp^k f)
      & = G (F (g \Pcomp^k f)) \\
      & = G (F (g) \Pcomp^k F (f)) \\
      & = G (F (g)) \Pcomp^k G (F (f)) \\
      & = (G \circ F) (g) \Pcomp^k (G \circ F) (f),
    \end{split}
  \]
  so $G \circ F$ preserves composition at dimension $k$.
\end{proof}

\begin{theorem}[Identity single-sorted functor]
  \label{thm:SingleSortedFunctor_id}
  \uses{def:SingleSortedFunctor}
  \lean{HigherCategoryTheory.SingleSorted.Functor.id}
  \leanok
  Let $\cat{C}$ be a single-sorted $\Index$-category. Then, the identity function $\id_C \colon C
  \to C$ is a single-sorted $\Index$-functor from $\cat{C}$ to itself.
\end{theorem}

\begin{proof}
  \leanok
  It trivially follows from reflexivity.
\end{proof}

\begin{definition}[Single-sorted finite functor]
  \label{def:SingleSortedNFunctor}
  \uses{def:SingleSortedNCategory, def:SingleSortedFunctor}
  \lean{HigherCategoryTheory.SingleSorted.NFunctor}
  \leanok
  Given a natural number $n \in \NN$, a \emph{single-sorted $n$-functor} is a single-sorted
  functor between single-sorted $n$-categories.
\end{definition}

\begin{definition}[Single-sorted omega functor]
  \label{def:SingleSortedOmegaFunctor}
  \uses{def:SingleSortedOmegaCategory, def:SingleSortedFunctor}
  \lean{HigherCategoryTheory.SingleSorted.OmegaFunctor}
  \leanok
  A \emph{single-sorted $\omega$-functor} is a single-sorted functor between single-sorted
  $\omega$-categories.
\end{definition}

\section{The category of single-sorted categories}

Lastly, we define the category whose objects are single-sorted categories and whose morphisms are
single-sorted functors.

\begin{definition}[Category of single-sorted categories]
  \label{def:SingleSortedCategory_category}
  \uses{
    def:SingleSortedCategory,
    def:SingleSortedFunctor,
    thm:SingleSortedFunctor_comp,
    thm:SingleSortedFunctor_id
  }
  \lean{HigherCategoryTheory.SingleSorted.Cat.category}
  \leanok
  The \emph{category of single-sorted $\Index$-categories}, denoted by $\Index\Cat$, is defined as
  follows:

  \begin{enumerate}
    \item \label{ax:ss_cat_obj} Its objects are single-sorted $\Index$-categories.
    \item \label{ax:ss_cat_mor} Its morphisms are single-sorted $\Index$-functors between
      single-sorted $\Index$-categories.
    \item \label{ax:ss_cat_comp} Composition of morphisms is given by composition of single-sorted
      functors, which is well defined by Theorem~\ref{thm:SingleSortedFunctor_comp}.
    \item \label{ax:ss_cat_id} The identity morphism for each object is given by the identity
      single-sorted functor, which is well defined by Theorem~\ref{thm:SingleSortedFunctor_id}.
  \end{enumerate}
\end{definition}

\begin{proof}
  \leanok
  The required category laws trivially follow from the corresponding properties of function
  composition.
\end{proof}

\begin{notation}
  Since single-sorted $n$-categories are just single-sorted $\Index$-categories where $\Index$ is
  the finite set $n$, it follows that we can define the category of single-sorted $n$-categories and
  single-sorted $n$-functors between them, denoted by $n\Cat$.
\end{notation}

\begin{definition}[Category of single-sorted omega categories]
  \label{def:SingleSortedOmegaCategory_category}
  \uses{
    def:SingleSortedOmegaCategory,
    def:SingleSortedOmegaFunctor,
    thm:SingleSortedFunctor_comp,
    thm:SingleSortedFunctor_id
  }
  \lean{HigherCategoryTheory.SingleSorted.OmegaCat.category}
  \leanok
  The \emph{category of single-sorted $\omega$-categories}, denoted by $\omega\Cat$, is defined
  analogously to Definition~\ref{def:SingleSortedCategory_category}, but using single-sorted
  $\omega$-categories and single-sorted $\omega$-functors instead.
\end{definition}

\begin{proof}
  \leanok
  As in Definition~\ref{def:SingleSortedCategory_category}, the required category laws trivially
  follow from the corresponding properties of function composition.
\end{proof}
